% Draft, 23-10-2026; n = 32768 and verified references 26-10-2026. Every number below is a measurement in ../README.md (same laptop, same code, the bench named beside it).
% Build: docker run --rm -v "$PWD":/w -w /w texlive/texlive:latest-small latexmk -pdf main.tex
\documentclass[11pt]{article}
\usepackage[T1]{fontenc}
\usepackage{lmodern}
\usepackage[a4paper,margin=2.4cm]{geometry}
\usepackage{amsmath,amssymb}
\usepackage{booktabs}
\usepackage{hyperref}
\usepackage{xcolor}
\hypersetup{colorlinks=true,linkcolor=blue!50!black,citecolor=blue!50!black,urlcolor=blue!50!black}

\title{Precision as a software choice:\\FP64, double-double, triple-double and beyond\\from the INT8 tensor cores of a laptop GPU}
\author{Shaban Ejupi}
\date{26 October 2026}

\begin{document}
\maketitle

\begin{abstract}
Consumer and laptop GPUs execute FP64 at 1/64 of their FP32 rate, while their INT8 tensor cores are several hundred times
faster than their FP64 units. On a laptop Blackwell GPU (RTX PRO 3000, sm\_120) we measure 232~TOPS of INT8
\texttt{mma.sync} against 0.449~TFLOPS of native FP64, a ratio of $517\times$. Building on the Ozaki scheme~II---exact integer
products of residues and reconstruction by the Chinese remainder theorem---we implement, in CUDA C and inline PTX, matrix
multiplication at FP64 (5.56~TFLOPS, $12.4\times$ the native peak and more accurate), double-double ($\sim$106 bits,
1.76~TFLOPS), triple-double ($\sim$159 bits, 1.58~TFLOPS) and arbitrary precision (212, 318 and 424 bits) from the same INT8
units. To our knowledge no GEMM beyond FP64 had been built this way. The higher precisions pay off in a real task: solving
$Ax=b$ with $\kappa(A)=10^{20}$ in double-double takes 0.855~s at $n=8192$, $499\times$ sooner than a
\texttt{\_\_float128} LU on all 24 CPU threads of the same laptop and with an error as small; at $\kappa=10^{36}$, where FP64
and double-double have no correct digit, a triple-double LU on the GPU is $32\times$ faster than the same arithmetic on 24 CPU
threads and $145\times$ faster than MPFR, and is the most accurate of the three. Counting CPU and GPU power, a double-double
solve uses $\sim$270$\times$ less energy than the quad-precision CPU solve. Leaving 8-bit moduli for primes below $2^{16}$ and $2^{23}$ split into bytes, an \emph{exact} product of integer matrices with entries of 1\,015 to 131\,063 bits costs work linear in the bits, is bit-exact against GMP, and runs $10$--$19\times$ faster than FLINT's multimodular \texttt{fmpz\_mat\_mul} on 24 CPU threads, with $15$--$18\times$ less energy. Turning the problem around---int8 in, residues out---the exact determinant of an int8 matrix comes $224\times$ sooner than FLINT at $n=4096$ and $355\times$ at $n=8192$ (98\,139 bits), with a proved bound on $|\det A|$ tight to one bit, and $n=16\,384$ (204\,480 bits) takes 66~s, $n=32\,768$ (425\,348 bits) 594~s within the 12~GB of the GPU. For entries of 32 to 256 bits, high-order lifting on the same modular GEMMs (the residue squared in RNS, the solution vector carried through INT8 GEMMs) and a proved bound for byte-plane entries give the determinant $236\times$ sooner than FLINT at $n=2048$ with 64-bit entries, $306\times$ with 128-bit entries and $265\times$ at $n=4096$ (276\,398 bits in 3.3~s against 868~s): the ratio grows with the entry size. The laptop is power-bound ($\sim$95~W) under these
kernels, which shapes what can and cannot be gained. On 15 of the most ill-conditioned real matrices of NIST Matrix Market ($n\le10\,000$, entries up to 154 bits) the certified exact solution takes 0.2--115~s, up to $\sim$900$\times$ sooner than SPEX; FP64 has no correct digit on 8 of them. The same units give correctly rounded MXFP4 and BF16 products: on the real weights of Kimi K3 (2.78 trillion parameters, the largest open-weights model) every output of 48 expert matrices is certified and identical for any batch, where BF16 tensor cores with FP32 accumulation round 51--83\% of them differently; at decode the exact kernel reads 94~GB/s, $18\times$ the laptop's NVMe, so reproducibility costs nothing when such a model streams from disk. On Nemotron~3 Ultra (560 billion parameters, the largest model whose training data are published) the GPU's own NVFP4 tensor-core instruction leaves 18--49\% of an expert's outputs not correctly rounded; one NVFP4 block is one INT8 MMA, and the correctly rounded product costs $4\times$ at instruction peak in prefill and nothing in decode.
\end{abstract}

\noindent\textbf{Përmbledhje (shqip).} Kartelat grafike të laptopëve e ekzekutojnë FP64 me 1/64 e shpejtësisë FP32, ndërsa
tensor core-t INT8 janë qindra herë më të shpejtë. Mbi një GPU Blackwell laptopi matëm 232~TOPS INT8 kundrejt 0,449~TFLOPS
FP64 ($517\times$). Me skemën Ozaki~II (prodhime të sakta të mbetjeve dhe teorema kineze e mbetjeve) ndërtuam shumëzim
matricash me saktësi FP64, double-double, triple-double dhe arbitrare (deri 424 bit) nga të njëjtat njësi INT8. Te një punë
reale---$Ax=b$ me $\kappa=10^{20}$---zgjidhja double-double mbaron $499\times$ më shpejt se LU \texttt{\_\_float128} mbi 24 fije
CPU, me gabim po aq të vogël, dhe me $\sim$270$\times$ më pak energji. Me prime nën $2^{16}$ e $2^{23}$ të ndara në bajte, shumëzimi i \emph{saktë} i matricave me numra të plotë 1\,015--131\,063 bitësh del bit për bit i saktë, $10$--$19\times$ më shpejt se FLINT dhe me $15$--$18\times$ më pak energji. Përcaktori i saktë i një matrice int8 del $224\times$ më shpejt se FLINT te $n=4096$ dhe $355\times$ te $n=8192$, me një kufi të provuar të ngushtë brenda një biti; te $n=32\,768$ (425\,348 bit) mbaron për 594~s brenda 12~GB të GPU-së. Për hyrje 32--256 bit, ngritja e rendit të lartë mbi të njëjtat GEMM modulare dhe një kufi i provuar për hyrje me plane bajtesh e japin përcaktorin $236\times$ më shpejt se FLINT te $n=2048$ me hyrje 64-bit, $306\times$ me 128-bit dhe $265\times$ te $n=4096$; raporti rritet me madhësinë e hyrjeve.

\section{Introduction}
The FP64 units of consumer and laptop GPUs are a small fraction of the chip: on our laptop Blackwell part a peak kernel
reaches 0.449~TFLOPS of FP64 against 28.2~TFLOPS of FP32 and 232~TOPS of INT8 tensor-core throughput (Table~\ref{tab:peak}).
Ozaki, Uchino and Imamura showed that FP64 matrix multiplication can be emulated on INT8 tensor cores (Ozaki scheme~II):
each input is scaled to integers, reduced modulo a set of pairwise coprime moduli, the residue matrices are multiplied
exactly in INT8 with INT32 accumulation, and the integer product is reconstructed by the Chinese remainder theorem (CRT)
before being scaled back~\cite{ozaki2}. FP64 by this route is not new; it is our baseline.

What we add is the observation, made concrete and measured, that the same machinery does not stop at FP64. The precision is
set by the number of moduli and by how many floating-point words each input and output carries; the expensive part, the
INT8 GEMMs, runs at the same speed whatever the precision. This turns precision into a software choice on hardware that has
no wide floating-point unit at all. We build and measure:
\begin{itemize}
\item FP64 GEMM at 16 moduli: 5.56~TFLOPS, $12.4\times$ the native FP64 peak, with smaller error than native FP64;
\item double-double ($\sim$106-bit) GEMM at 32 moduli: 1.76~TFLOPS, $3.9\times$ the native FP64 peak with twice the bits;
\item triple-double ($\sim$159-bit) GEMM at 49 moduli---after which every prime below 256 divides one of them, so no further
8-bit modulus can join---at 1.58~TFLOPS;
\item arbitrary precision beyond that from balanced 160-bit slices: 212, 318 and 424 bits;
\item dense linear solves at condition numbers where FP64 and double-double fail, all on the GPU.
\item exact determinants of int8 matrices up to $n=32\,768$, and of matrices with 32--256-bit entries by high-order lifting on tensor cores, 63--306$\times$ faster than FLINT.
\end{itemize}
Everything runs on one laptop, and all competitors were measured on the same machine.

\begin{table}[t]\centering
\caption{Peak throughput of the laptop GPU (\texttt{bench/peak.cu}).}\label{tab:peak}
\begin{tabular}{lr}\toprule
Unit & Throughput \\\midrule
FP64 FMA, native & 0.449 TFLOPS \\
FP32 FMA, native & 28.2 TFLOPS \\
INT8 \texttt{mma.sync.m16n8k32} (PTX) & 232 TOPS \\
\textbf{INT8 / FP64} & $\mathbf{517\times}$ \\\bottomrule
\end{tabular}
\end{table}

\section{Method}
\paragraph{Ozaki scheme II.} For $C=AB$ the rows of $A$ and columns of $B$ are scaled by powers of two to integers
$A',B'$ whose product cannot exceed $M/2$ in magnitude, where $M=\prod_\ell m_\ell$ is the product of $L$ pairwise coprime
moduli $m_\ell\le 256$. For each $\ell$ the residues $A'\bmod m_\ell$ and $B'\bmod m_\ell$ fit in INT8, their product is
computed exactly by the tensor cores with INT32 accumulation, and reduced mod $m_\ell$. The CRT then gives
$A'B'$ exactly as a signed integer of up to $\log_2 M$ bits, which is scaled back. The precision of the result is
limited only by how many bits of each input survive the scaling, i.e.\ by $\log_2 M$.

\paragraph{Wider inputs and outputs.} For double-double and triple-double inputs, each of the two or three parts is split
exactly into 53-bit integer pieces before the residues are formed; like pieces are summed and cut once into 16-bit chunks,
so the residue of a wide integer costs 32-bit multiply-adds per modulus instead of 64-bit reductions (triple-double residues
57~$\to$~19~ms at $n=4096$). On the output side the integer is never converted through floating-point arithmetic: the CRT is
evaluated as a direct sum $x=\sum_\ell r_\ell W_\ell - qM$ over 32-bit limbs on the integer units, with $q$ from 40-bit
fixed-point fractions, and the result is peeled off the integer as two or three doubles (round down, round down, round to
nearest, each exact from the integer), then normalized. This replaced a Garner reconstruction followed by a double-double
Horner evaluation, which took 41.5 of 47~ms of a $4096^2$ double-double GEMM because FP64 runs at 1/64 rate on this chip;
the direct sum costs 0.68~ns per output element instead of 2.8 and is nearly independent of $K$. It was checked on the host
against $10^5$ random 200-bit integers with no error.

\paragraph{Moduli.} FP64 uses 16 moduli, double-double 32, triple-double 49 (the 32 plus the 17 primes 103\ldots 29 that are
not yet factors; $M\approx 2^{342}$). With these, every prime below 256 divides one of the moduli, so no further 8-bit modulus
coprime to all of them exists and $M$ cannot grow within one GEMM. Beyond $\sim$160-bit inputs we cut each input into balanced 160-bit slices and run one 49-modulus GEMM per
slice pair that reaches the result, accumulating exactly in wide integers and peeling $Q$ doubles.

\section{Accuracy}
Errors are the maximum componentwise $|C-C_{\mathrm{ref}}|/(|A||B|)$ in units of the target precision's unit roundoff,
against higher-precision CPU references (80-bit for FP64, \texttt{\_\_float128} for double-double, an exact 600-bit MPFR sum for
triple-double, an exact 2400-bit sum beyond). Inputs have magnitudes spread over $2^{\pm 0}$ or $2^{\pm 8}$.

\begin{table}[h]\centering
\caption{Error in units of $2^{-53}$ (FP64, $M=N=256$) and $2^{-106}$ (double-double, $M=N=128$).}
\begin{tabular}{rrrrr}\toprule
$K$ & magnitudes & native FP64 & 14 moduli & \textbf{16 moduli} \\\midrule
512 & $2^{\pm0}$ & 4.71 & 10.9 & \textbf{0.99} \\
512 & $2^{\pm8}$ & 24.1 & 395 & \textbf{3.1} \\
4096 & $2^{\pm0}$ & 4.15 & 10.5 & \textbf{0.25} \\
4096 & $2^{\pm8}$ & 25.1 & 203 & \textbf{1.87} \\\midrule
$K$ & magnitudes & 24 moduli & 28 moduli & \textbf{32 moduli} \\\midrule
256 & $2^{\pm0}$ & $6.2\cdot10^5$ & 27.1 & \textbf{1.09} \\
256 & $2^{\pm8}$ & $2.0\cdot10^7$ & 1160 & \textbf{46.9} \\
1024 & $2^{\pm0}$ & $4.7\cdot10^5$ & 31 & \textbf{0.49} \\
1024 & $2^{\pm8}$ & $1.1\cdot10^7$ & 713 & \textbf{12.4} \\\bottomrule
\end{tabular}
\end{table}

\begin{table}[h]\centering
\caption{Triple-double error in units of $2^{-159}$, $M=N=128$, against an exact MPFR sum.}
\begin{tabular}{rrrrr}\toprule
$K$ & magnitudes & 41 moduli & 45 moduli & \textbf{49 moduli} \\\midrule
256 & $2^{\pm0}$ & $1.7\cdot10^4$ & 5.2 & \textbf{0.24} \\
256 & $2^{\pm8}$ & $1.2\cdot10^6$ & 343 & \textbf{10.7} \\
1024 & $2^{\pm0}$ & $1.8\cdot10^4$ & 4.6 & \textbf{0.13} \\
1024 & $2^{\pm8}$ & $4.3\cdot10^5$ & 110 & \textbf{2.1} \\\bottomrule
\end{tabular}
\end{table}
Beyond triple-double, against an exact 2400-bit sum at $K=256$: 0.13~ulp at 212 bits, 0.25 at 318 bits, 0.08 at 424 bits.

\section{Speed}
\begin{table}[h]\centering
\caption{GEMM speed on the laptop. CPU numbers with the CPU otherwise idle, measured at $n=256$/512 and scaled by $n^3$.}\label{tab:speed}
\small
\setlength{\tabcolsep}{4pt}
\begin{tabular}{rrrrll}\toprule
bits & $n$ & GPU time & GPU TFLOPS & best CPU code, 24 threads & MPFR, 24 threads \\\midrule
53 (FP64) & 4096 & --- & 5.56 ($12.4\times$ native) & --- & --- \\
106 (dd) & 4096 & --- & 1.76 & --- & --- \\
159 (td) & 4096 & 87 ms & 1.58 & td AVX $\sim$10.7 s ($123\times$) & $\sim$208 s ($2390\times$) \\
212 & 2048 & 36 ms & 0.47 & QD \texttt{qd\_real} $\sim$16.2 s ($445\times$) & $\sim$35 s ($960\times$) \\
318 & 2048 & 52 ms & 0.33 & --- & $\sim$37 s ($713\times$) \\
424 & 2048 & 80 ms & 0.22 & --- & $\sim$41 s ($514\times$) \\\bottomrule
\end{tabular}
\end{table}
At 16 moduli the time of a $4096^3$ FP64 GEMM splits into scales 1.0~ms, residues 3.0~ms, the 16 INT8 GEMMs 16.8~ms and the
CRT 3.6~ms. The INT8 GEMM reaches 157~TOPS, 67\% of the measured peak, with a three-stage \texttt{cp.async} pipeline; a faster
kernel did not move end-to-end times, because with real data the laptop is power-bound (Section~\ref{sec:energy}).

\section{A real task: ill-conditioned linear systems}
We solve $Ax=b$ with $A=Q_1\,\mathrm{diag}(\sigma)\,Q_2$, $Q_{1,2}$ orthogonal (a DST-I, the second permuted and sign-flipped),
$\sigma$ spaced from 1 to $10^{-20}$ or $10^{-36}$; such matrices cannot be fixed by scaling. $b=Ax_{\mathrm{true}}$ is summed
exactly. The error is $\max|x-x_{\mathrm{ref}}|/\max|x_{\mathrm{ref}}|$.

\paragraph{Double-double, $\kappa=10^{20}$.} $\kappa u$ is $10^4$ for FP64, so FP64 LU has no correct digit and textbook
iterative refinement diverges (it needs $\kappa u<1$; at $\kappa=10^{12}$ the same refinement converges to $3\cdot10^{-20}$, so
the baseline is right). The GPU solver keeps everything on the device: a recursive panel with cooperative 32-column leaves (a
thread per row holding the row in registers, rows renamed instead of swapped, one grid synchronization per column), halves
joined by double-double GEMMs on the tensor cores, $L_{11}^{-1}$ assembled from the halves' inverses, left-looking solves with
inverted diagonal blocks, and one refinement step in double-double.

\begin{table}[h]\centering
\caption{Double-double LU on the GPU vs \texttt{\_\_float128} LU on 24 CPU threads, $\kappa=10^{20}$, nb $=1024$.}
\begin{tabular}{rrrrl}\toprule
$n$ & dd LU, GPU & \texttt{\_\_float128} LU & speedup & error vs $x_{\mathrm{true}}$ (ours / quad) \\\midrule
2048 & 0.113 s & 7.65 s & $67\times$ & $1.6\cdot10^{-11}$ / $2.4\cdot10^{-11}$ \\
4096 & 0.283 s & 55.6 s & $\sim197\times$ & --- \\
6144 & 0.572 s & 200.8 s & $351\times$ & $3.4\cdot10^{-10}$ / $2.9\cdot10^{-10}$ \\
8192 & 0.956 s & 476.9 s & $\mathbf{499\times}$ & $2.1\cdot10^{-11}$ / $2.3\cdot10^{-11}$ \\\bottomrule
\end{tabular}
\end{table}
The 0.956~s at $n=8192$ was measured in the same run as the quad reference, after minutes of all-core load; alone on a cool
GPU the solve takes 0.855~s. Of those, the panel takes 0.45~s (its leaves 0.22~s, bound by the FP64 units at
$\sim$10~GFLOPS per SM), the trailing GEMMs 0.15~s, solves and refinement 0.09~s, $U_{12}$ 0.04~s.

\paragraph{Triple-double, $\kappa=10^{36}$.} Here $\kappa u$ is $10^{20}$ for FP64 and $10^4$ for double-double: neither has a
correct digit. The same solver structure in three parts per number, with 16-column leaves:

\begin{table}[h]\centering
\caption{Triple-double LU, $\kappa=10^{36}$; GPU speedups in parentheses; CPU otherwise idle.}
\begin{tabular}{rllll}\toprule
$n$ & \textbf{td LU, GPU} & td LU, 24 CPU threads & MPFR LU, 160 bits & FP64 LU \\\midrule
1024 & \textbf{0.082 s}, $1.3\cdot10^{-14}$ & 0.10 s, $3.6\cdot10^{-14}$ ($1.2\times$) & 1.32 s, $3.3\cdot10^{-14}$ ($16\times$) & 128 \\
2048 & \textbf{0.212 s}, $2.1\cdot10^{-14}$ & 1.04 s, $7.6\cdot10^{-14}$ ($5\times$) & 10.2 s, $5.6\cdot10^{-14}$ ($48\times$) & 103 \\
4096 & \textbf{0.563 s}, $2.2\cdot10^{-14}$ & 18.0 s, $1.2\cdot10^{-13}$ ($32\times$) & 81.9 s, $9.6\cdot10^{-14}$ ($145\times$) & 86 \\\bottomrule
\end{tabular}
\end{table}
At $n\le1024$ a well-vectorized triple-double LU on 24 cores keeps up; the gap opens with $n$ as the work moves into the
tensor-core GEMMs. At $n=4096$ the panels take 0.38~s of the 0.563~s (leaves 0.14~s): the panel, not the GEMM, is now the
limit.

\section{Thousands of bits, exactly: linear in the bits}\label{sec:big}
The 49 moduli $\le 256$ cap one Ozaki-II product at $\sim$160-bit inputs, and slicing beyond that costs a number of GEMMs that
grows with the square of the slices. For wide integers we leave 8-bit moduli: the moduli are primes $p<65\,280$, and each
balanced residue ($|r|<p/2$) is split into two signed bytes, $r=256h+\ell$. Then $AB \bmod p$ is four exact INT8 GEMMs
($hh$, $h\ell$, $\ell h$, $\ell\ell$) recombined in 64-bit integers, so \emph{the GEMM work grows linearly with the number of
bits}. Above $\sim$60\,000 bits of the CRT modulus the residues take three bytes and the primes run up to $2^{23}$ (nine
products per prime, $\sim$23 bits each). The two conversions are INT8 GEMMs as well, as conversions to and from a residue
number system are matrix products: the residues are the balanced byte digits of each entry times the bytes of
$2^{8j} \bmod p$; for the CRT, $t_\ell = r_\ell\,(M/p_\ell)^{-1} \bmod p_\ell$ (bytes) times the bytes of $M/p_\ell$ gives the digits
of $x+qM$, with $q=\operatorname{round}(\sum_\ell t_\ell/p_\ell)$ from doubles, and $qM$ is taken off while the digits are
carried. Everything runs on the tensor cores (\texttt{src/ozaki\_big.cuh}, \texttt{bench/big\_test.cu}).

Every product is exact: 256 sampled elements per run agree bit for bit with GMP. The baselines compute the same exact product
on all 24 CPU threads: GMP inner products (\texttt{mpz\_addmul}), and FLINT's \texttt{fmpz\_mat\_mul}, which is itself
multimodular with BLAS-style kernels and is the strongest baseline we know of.
\begin{center}\begin{tabular}{rrrrrrr}\toprule
entry bits & $n$ & primes (bytes) & GPU & GMP & FLINT & GPU vs FLINT \\\midrule
1\,015 & 1024 & 192 (2) & 71 ms & 23.0 s & 1.31 s & $18\times$ \\
4\,087 & 512 & 576 (2) & 62 ms & 16.4 s & 1.07 s & $17\times$ \\
16\,375 & 256 & 2\,112 (2) & 90 ms & 14.0 s & 1.69 s & $19\times$ \\
4\,087 & 1024 & 640 (2) & 341 ms & 125 s & 6.10 s & $18\times$ \\
16\,375 & 512 & 2\,112 (2) & 363 ms & 104 s & 6.33 s & $17\times$ \\
32\,759 & 256 & 4\,352 (2) & 316 ms & 36.1 s & 4.56 s & $14\times$ \\
65\,527 & 128 & 5\,760 (3) & 313 ms & 10.6 s & 3.45 s & $11\times$ \\
131\,063 & 128 & 11\,520 (3) & 983 ms & 29.6 s & 9.76 s & $10\times$ \\\bottomrule
\end{tabular}\end{center}
GMP times are measured on a corner of the outputs ($96^2$, or $48^2$ above 8\,000 bits) and scaled; FLINT runs the full
product. The largest entries have $\sim$39\,450 decimal digits. Energy per exact product (RAPL package + \texttt{nvidia-smi},
both chips counted in both columns): 1\,015 bits, $n=1024$: 8.6~J on the GPU path against 131~J for FLINT ($15\times$);
16\,375 bits, $n=256$: 7.9~J against 142~J ($18\times$).

Doubling $n$ at fixed width does not raise the ratio either (16\,375 bits: $19\times$ at $n=256$, $17\times$ at $n=512$): even at $n=512$ the conversion GEMMs ($\sim 1.9\cdot10^{13}$ INT8 operations) dwarf the 8\,704 product GEMMs ($\sim 2.3\cdot10^{12}$). The ratio to FLINT falls slowly beyond $\sim$16\,000 bits ($19\times\to10\times$) because the conversions are $n^2P^2$ work
for $P$-bit entries while the products are $n^3P$: at $n=128$ the conversions dominate. The way past that crossover is a
subquadratic conversion (a two-level residue system), which we leave for the next version. For scale: one part in a googol
($10^{-100}$) is $\sim$333 bits, already inside the 424-bit format of the previous section.


\section{Changing the problem: exact determinants and exact solutions}\label{sec:det}
The ratio of the previous section falls with the bits because large numbers must enter and leave the residue system. Many
exact computations produce huge numbers from small inputs: the determinant of an $n\times n$ int8 matrix has $\sim$0.72 bits per
row less than Hadamard's bound, about 47\,000 bits at $n=4096$, and the rational solution of $Ax=b$ has $n$ numerators of that
size. There the conversions vanish: the input is int8, each prime returns one residue (or a vector of digits), and the only
big-number step is on the host (\texttt{src/ozaki\_det.cuh}, \texttt{bench/det.cu}).

\paragraph{LU modulo thousands of primes.} For primes $p<65\,280$ and residues $X, W$, with $X=X_\ell+256X_h$ in balanced bytes and
$V=256W \bmod p$, $XW \equiv [X_\ell\,X_h][W_\ell;V_\ell] + 256\,[X_\ell\,X_h][W_h;V_h] \pmod p$: two INT8 GEMMs of depth $2K$,
reduced once in the epilogue. A recursive LU (Toledo) without pivoting runs for a batch of primes at once, every $O(n^3)$ step
such a product. Its 64-column leaves invert their $64\times64$ diagonal blocks in shared memory. A prime with a vanishing
leading minor (probability $\sim n/p$) is dropped. Structured matrices, whose leading minors vanish over $\mathbb Z$, are
preconditioned as $UAL$ with random unit triangular $U, L$ (Kaltofen--Saunders) whose entries in $[0,127]$ are the same residues for every prime.

\paragraph{A divisor by $p$-adic lifting.} (The divisor-then-primes hybrid is classical on CPUs~\cite{pan}; FLINT uses it.) $A^{-1}\bmod p$ comes from the LU of $[A\;I]$. Dixon lifting runs on the GPU, one
16-bit and one int8 matrix--vector product per step: $x_i=A^{-1}r_i \bmod p$, $r_{i+1}=(r_i-Ax_i)/p$. Rational reconstruction of a
few entries of $x=A^{-1}b$ gives $d$, which divides $\det A$ and for random $b$ is almost always the largest invariant factor.
Then $\det A/d$ needs only enough primes to exceed a bound on $|\det A|/d$.

\paragraph{A proved bound tight to one bit.} Hadamard's bound leaves $\sim0.72n$ bits of slack on a random matrix (5\,900 at
$n=8192$), i.e.\ hundreds of primes for $\det A/d$. We replace it by a certificate in the spirit of Ogita's verified
determinants~\cite{ogita}. FP32 LU with partial pivoting gives $PA\approx LU$. $X_L=\lfloor2^sL^{-1}\rceil$ (unit diagonal,
so exactly $2^s$) and $X_U=\lfloor2^sU^{-1}\operatorname{diag}(u)\rceil$ are integer triangular matrices, so their determinants
are products of their diagonals. $B=X_L\,PA\,X_U$ is computed \emph{exactly} (the same modular INT8 GEMMs, int64 operands
reduced while packing, Garner to 128 bits), and Hadamard's inequality on $B$, accumulated with upward rounding, bounds
$|\det A| = |\det B|/(2^{sn}\,|\det X_U|)$. Floating point only decides how close $B$ is to $2^{2s}\operatorname{diag}(u)$,
that is how tight the bound is, never whether it holds. Measured: 47\,021 bits against $|\det A|$ of 47\,021 bits at $n=4096$,
98\,140 against 98\,139 at $n=8192$. $\det A/d$ then needs 6 primes instead of 224 and 496.

\paragraph{Results.} Random int8 matrices, FLINT~3.0.1 on all 24 CPU threads; every result bit-exact against FLINT.
FLINT's \texttt{fmpz\_mat\_det} is already its accelerated path (a Dixon divisor, then primes). Its plain multimodular
\texttt{fmpz\_mat\_det\_modular} takes 347~s at $n=2048$, $177\times$ our all-primes LU with the same algorithm.
\begin{center}\small\begin{tabular}{rrrrrr}\toprule
$n$ & $|\det A|$ bits & GPU, all primes & GPU, lifting + bound & FLINT & vs FLINT \\\midrule
1\,024 & 10\,736 & 0.25 s & 0.073 s & 2.89 s & $40\times$ \\
2\,048 & 22\,485 & 1.96 s & 0.293 s & 24.9 s & $85\times$ \\
4\,096 & 47\,021 & 21.3 s & 1.15 s & 257 s & $224\times$ \\
8\,192 & 98\,139 & --- & 8.05 s & 2\,861 s & $355\times$ \\
16\,384 & 204\,480 & --- & 65.9 s & --- & \\
32\,768 & 425\,348 & --- & 594 s & --- & \\\bottomrule
\end{tabular}\end{center}
At $n=8192$ FLINT ran while some GPU runs shared the CPU; the leading digits agree, and $n=8192, 16\,384, 32\,768$ are also checked
modulo four primes used nowhere else. The exact rational solution of $Ax=b$ (same lifting; all $n$ entries back to
integers by Horner in halves on 24 threads; one common denominator) is bit-exact against FLINT's
\texttt{fmpz\_mat\_solve\_dixon\_den}: 0.064~s against 2.33~s at $n=1024$ ($36\times$), 0.145~s against 18.9~s at $n=2048$
($130\times$), 1.62~s against 145~s at $n=4096$ ($90\times$, 192.6~Mbit of numerators), and 7.5~s at $n=8192$ (804~Mbit).
Energy per determinant at $n=2048$, both chips counted: $\sim$40~J on the GPU path (process start-up included; $\sim$23~J for the
computation alone) against $\sim$2\,030~J for FLINT, $\sim$50$\times$ less. At $n=8192$ the lifting is the cost: 13\,013 steps,
each reading $A^{-1}$ (16 bits per entry) and $A$, at $\sim$100\% of the read bandwidth we measure on this GPU ($\sim$410~GB/s).
That limit is physical for digit-by-digit lifting; moving past it needs lifting built on matrix products (high-order lifting,
the DoublePlusOneLift of Pauderis and Storjohann~\cite{ps}), the next step.

\paragraph{$n=32\,768$ in 12~GB.} One modular product packs $(2n+128)\cdot 2n$ bytes three times (13~GB), and the bound holds
three float and two int64 $n\times n$ matrices (30~GB). Products deeper than 2048 now run in $K$-chunks (a third epilogue
mode, $R \mathrel{+}= XW$); the bound builds $X_L$, $X_U$ in column panels (triangular solves against identity panels) kept in
host memory, forms $T=X_L\,PA \bmod q$ in row panels and $B=T X_U$ in column panels as 16-bit residues on the host, and sums
the rows in panels, the $L_1$ norms exactly in 128 bits; it runs on its own stream beside the lifting. Peak use: $\sim$11~GB
of VRAM and 29~GB of host memory. The determinant (425\,348 bits) takes 594~s: $A^{-1}$ 15~s, 56\,145 lifting steps 497~s at
the bandwidth limit, $\det A/d$ 39 primes 71~s; the bound, 244 bits above $|\det A|$ at this size (FP32 LU; a larger scale
$2^s$ does not change it), is hidden under the lifting. Each piece was checked bit-exact against FLINT at $n\le 4096$ with the
chunks and panels forced small.
\paragraph{Large entries.} With entries of $L$ bytes stored as byte planes ($A=\sum_j 256^j A_j$), the residues come by Horner and the LU is
unchanged; the lifting reads the $L$ planes per step and keeps the residual as $L+4$ base-256 digits per row, divided by $p$
exactly by one lane. Lifting modulo $M=p_0\cdots p_{S-1}$ ($S\approx L/2$ primes at once: $x_i$ by CRT from $S$ residues,
$A x_i$ on byte-aligned digits, division by each $p_t$ in turn) gains $16S$ bits per read of the $L$ planes and cuts the
lifting 1.3--2.9$\times$. Bit-exact against FLINT: $n=1024$ with 64-bit entries in 0.247~s against 12.4~s ($50\times$), with
128-bit entries 0.692~s against 31.0~s ($45\times$); $n=2048$ with 32-bit entries 0.563~s against 46.7~s ($83\times$), with
64-bit entries (a 137\,173-bit determinant) 2.03~s against 106.9~s ($53\times$). GPU times come from runs without FLINT beside
them: CPU and GPU share one power budget, and right after 24 FLINT threads the GPU is up to $2\times$ slower. Each step still
reads all $L$ planes, so the cost grows faster than the bits.

\paragraph{High-order lifting.} Digit-by-digit lifting reads $A$ and $A^{-1}$ once per $X$-digit; high-order lifting
(Storjohann's sparse inverse expansion~\cite{stor05}, in row form $yA=b$) squares the residue instead, so the precision
doubles per level. With $X=p_0\cdots p_{s-1}\ge 6.25\,n^2\|A\|$ (distinct 16-bit primes, $B_0=A^{-1}\bmod X$ from the same
batched LU), $R_0=(I-AB_0)/X$ and, per level, $\bar R=R_i^2$, $M_i=B_0\bar R \bmod X$, $R_{i+1}=(\bar R-AM_i)/X$ exactly, with
$\|R_i\|\le 1.25\,n\|A\|$; then $c_{i+1}=c_i+X^{m_i}c_iR_i+X^{2m_i}bM_i$, $m_{i+1}=2m_i+1$, gives $c_i\equiv bA^{-1}\pmod{X^{m_i}}$.
The residue chain runs in RNS, modulo the $s$ primes of $X$ and $w$ more primes ($\prod q>2.5\,n\|A\|$): $M_i$ by
Garner from its residues, $R_{i+1}\bmod q=(\bar R-AM_i)X^{-1}$, its exact value by Garner over the $q$ (checked against its
bound) --- $2(s+w)$ modular GEMMs per level and no big number anywhere. All matrices are kept transposed so that the vector
chain reads $K$-contiguous rows. That chain is itself a GEMM: $c_i$ is stored as balanced digits in a periodic mixed radix
(position $e$ has radix $p_{e\bmod s}$), so $X^m$ is a shift by $sm$ positions and the Garner digits of $M_i$ are already
digits of it; rows are the digits as $[\mathrm{lo}\,|\,\mathrm{hi}]$ bytes, columns the byte planes of $R_i$ as
$[Z_v\,|\,Z_{v-1}]$, so column $v$ collects $256^v$; the planes are combined in parallel and the carries run down the
positions one thread per entry. Around it, the proved bound now takes byte planes (a float LU of $A/256^{L-1}$ only decides
tightness; $B=X_L\,PA\,X_U$ is exact from the planes, Garner over up to 16 primes with $|B_{ij}|$ rounded up from the
mixed-radix digits): at $n=2048$ with 64-bit entries it is 1 bit above $|\det A|$, and $\det A/d$ needs 6 primes instead of
128. The rational reconstruction is FLINT's half-gcd \texttt{\_fmpq\_reconstruct\_fmpz\_2} (0.49~s $\to$ 0.047~s there).
Bit-exact against FLINT in every case of Table~\ref{tab:hol}: the lifting is 2.6--4.3$\times$ faster than $S$-prime Dixon and
grows like the bits, not their square, so the ratio to FLINT now rises with the entry size---$63\times$ at $n=512$ with
256-bit entries, $236\times$ at $n=2048$ with 64-bit and $306\times$ with 128-bit entries, $265\times$ at $n=4096$ with
64-bit entries (a 276\,398-bit determinant in 3.27~s against 868~s). At that size the residue chain is 2.08~s of the 2.71~s
lifting (247 modular GEMMs of $n^3$, power-bound), the vector chain 0.44~s.
\begin{table}[t]\centering
\caption{Exact determinants of $n\times n$ matrices with random $b$-bit entries: total GPU time with $S$-prime Dixon lifting
and with high-order lifting (both with the proved bound for byte planes and FLINT's reconstruction), against FLINT
\texttt{fmpz\_mat\_det} on 24 threads. GPU runs without FLINT beside them; all bit-exact.}\label{tab:hol}
\begin{tabular}{rrrrrrr}\toprule
$n$ & $b$ & $|\det A|$ bits & Dixon & HOL & FLINT & vs FLINT \\\midrule
256 & 256 & 65\,914 & 0.139~s & 0.063~s & 1.08~s & $17\times$ \\
512 & 256 & 132\,082 & 0.399~s & 0.139~s & 8.82~s & $63\times$ \\
1024 & 64 & 68\,073 & 0.177~s & 0.126~s & 12.39~s & $98\times$ \\
1024 & 128 & 133\,615 & 0.434~s & 0.219~s & 30.98~s & $141\times$ \\
2048 & 32 & 71\,638 & 0.358~s & 0.288~s & 46.70~s & $162\times$ \\
2048 & 64 & 137\,173 & 1.091~s & 0.453~s & 106.89~s & $236\times$ \\
2048 & 128 & 268\,249 & 2.946~s & 0.856~s & 261.65~s & $306\times$ \\
4096 & 32 & 145\,330 & 3.425~s & 1.934~s & 416.46~s & $215\times$ \\
4096 & 64 & 276\,398 & 7.443~s & 3.273~s & 867.68~s & $265\times$ \\\bottomrule
\end{tabular}
\end{table}
Moduli below 256 for $X$ (one GEMM per modulus, as in Ozaki-II) are not an option here: without pivoting the LU modulo $p$
fails with probability $\approx 1-e^{-n/p}$.

\section{Real systems: how wrong FP64 is, and the exact answer}\label{sec:real}
The residue chain of the lifting above now runs over 8-bit moduli wherever B$_0$ is not involved: $Z^2$ exactly in an RNS of
23--42 odd moduli below 256 (one INT8 GEMM each; on random data a 16-bit modular product costs $2.8\times$ more per bit of
modulus at $n=4096$), a base extension by approximate CRT to the 16-bit primes of $X$, and $AM$ modulo the first 10--19 moduli
only. At $4096\times64$ bits this takes the determinant from 3.22 to 2.28~s ($380\times$ FLINT), with the same result.

We then took the 13 most ill-conditioned square real matrices with $900\le n\le 8192$ of NIST Matrix Market---fluid dynamics
(FIDAP, LNS), magnetohydrodynamics (MHD), the Quebec Hydro network, crystal growth, oil reservoirs, a chemical
plant---and solved $Ax=b$ \emph{exactly}. Every double is $m\,2^e$, so scaling row $i$ by $2^{k_i}$ gives integer entries (up
to 154 bits) and the same $x$, with no rounding anywhere; rows are ordered by FP64 partial pivoting, since the GPU LU does not
pivot. High-order lifting gives all $n$ entries, a common denominator follows from a few half-gcd reconstructions, the
reconstruction is unique by the Hadamard and Cramer bounds, and $yM=c$ is checked modulo two further primes. On the same
integers we ran SPEX Left LU~\cite{spex} (SuiteSparse 7.6, exact sparse REF LU, one CPU thread); both programs print a digest
$\sum_j x_j r^j \bmod 2^{61}-1$. FP64 is LAPACK \texttt{dgesv} and SuperLU on the original doubles.

\begin{table}[h]\centering\footnotesize\setlength{\tabcolsep}{3pt}
\begin{tabular}{lrrrrrrcrr}
\hline
matrix & $n$ & $\kappa$ (NIST) & bits & GPU s & SPEX s & & same $x$ & \multicolumn{2}{c}{FP64 digits (entries with none)}\\
 & & & & & & & & LAPACK & SuperLU \\ \hline
west0989 & 989 & 5.7$\cdot 10^{12}$ & 75 & 0.17 & 1.5 & 9$\times$ & yes & 11.7 (0) & 11.8 (0) \\
qh1484 & 1484 & 5.6$\cdot 10^{17}$ & 96 & 0.55 & 12.4 & 23$\times$ & yes & 11.2 (0) & 12.2 (0) \\
fidap032 & 1159 & 5.2$\cdot 10^{20}$ & 60 & 0.28 & 54.5 & 193$\times$ & yes & 0.0 (568) & 0.0 (572) \\
fidap006 & 1651 & 3.5$\cdot 10^{21}$ & 154 & 1.08 & 963.7 & 895$\times$ & yes & 0.0 (887) & 0.0 (693) \\
cry2500 & 2500 & 4.4$\cdot 10^{17}$ & 70 & 1.41 & 699.3 & 497$\times$ & yes & 6.4 (0) & 5.1 (0) \\
fidapm07 & 2065 & 6.8$\cdot 10^{19}$ & 118 & 1.19 & $>$7200 &  & -- & 0.0 (430) & 0.0 (290) \\
fidap026 & 2163 & 4.7$\cdot 10^{18}$ & 153 & 2.12 & $>$7200 &  & -- & 0.0 (1568) & 0.0 (1530) \\
mhd3200a & 3200 & 6.5$\cdot 10^{39}$ & 130 & 4.02 & 1290.8 & 321$\times$ & yes & 0.0 (2876) & 0.0 (2768) \\
fidap014 & 3251 & 2.6$\cdot 10^{16}$ & 124 & 4.74 & $>$7200 &  & -- & 7.6 (0) & 8.0 (0) \\
lns\_3937 & 3937 & 1.0$\cdot 10^{17}$ & 79 & 4.31 & -- &  & -- & 10.1 (2) & 7.4 (1) \\
fidap012 & 3973 & 2.2$\cdot 10^{20}$ & 119 & 7.44 & -- &  & -- & 0.0 (483) & 0.0 (482) \\
mhd4800a & 4800 & 3.2$\cdot 10^{47}$ & 131 & 34.39 & -- &  & -- & 0.0 (3923) & 0.0 (4754) \\
sherman3 & 5005 & 6.9$\cdot 10^{16}$ & 102 & 14.25 & -- &  & -- & 11.8 (0) & 12.4 (0) \\
fidapm15 & 9287 & 1.36$\cdot 10^{20}$ & 63 & 78.66 & -- &  & -- & 0.0 (2251) & 0.0 (5260) \\
cry10000 & 10000 & 6.9$\cdot 10^{20}$ & 76 & 115.02 & -- &  & -- & 4.2 (66) & 1.6 (923) \\
\hline
\end{tabular}
\caption{Exact solutions of real systems on the laptop (GPU: $A^{-1}$, lifting, reconstruction and all numerators). SPEX
ran afterwards on the CPU alone, limit 2~h. $yM=c$ held modulo both check primes in every row.}\label{tab:real}
\end{table}

On 7 of the 13 systems FP64 has not one correct digit, LAPACK and SuperLU alike---on mhd3200a and mhd4800a the normwise error
is $10^8$ and 82--99\,\% of the entries have no correct digit---and NIST's own (FP64) condition estimate does not say which:
qh1484 at $5.6\cdot10^{17}$ keeps 11 digits. The certified exact solution takes 0.2--34~s; SPEX returns the same $x$, bit for
bit, where it finishes within the limit, 9 to $\sim$900$\times$ later.

\section{The largest open model: correctly rounded MXFP4 and BF16 products}\label{sec:llm}
\sloppy
Kimi K3 (Moonshot AI; read from its Hugging Face model card and safetensors headers: 2\,779\,931\,837\,184 parameters, 1.56~TB)
is the largest open-weights model. Its routed experts (2.72~T parameters; 93 layers $\times$ 896 experts, 16 per token) are
stored in MXFP4: an E2M1 value per weight, one E8M0 scale per 32. A group of 32 is the depth of one \texttt{mma.m16n8k32}
INT8 step and $2\cdot$E2M1 is an integer $\le 12$, so MXFP4$\times$BF16 products can be computed \emph{exactly} on INT8 tensor
cores and rounded once (\texttt{llm/mxexact.cu}, \texttt{llm/exact.cuh}). Weights: $q' = 2\,\mathrm{E2M1}\cdot 2^{(s-s_{\min})\bmod 4}$ is
an int8 ($\le 96$), one scale per row and band of four binades (every row of the 48 K3 matrices tested spans $\le 3$ binades:
one band). Activations: each token is aligned to its largest exponent and cut into $S$ int8 digits ($D_c = 7S-8$ bits below
the top), smaller elements form further tiers. Every (band, tier, digit) pair is a plain INT8 GEMM whose int32 sums cannot
overflow ($|P| \le K\cdot 96\cdot 128 < 2^{26.4}$); the few terms of an output are added in a 128-bit integer and rounded once,
with a certificate when a term falls below the 94-bit window (it never failed). BF16$\times$BF16 is the same with both sides
in digits. A correctly rounded result cannot depend on the summation order, the tiling, the batch or the machine.

On the real weights of layer 5 (experts 0--15, all three matrices, fetched by HTTP range requests; Student-t activations
scaled by the layer's real norm weight), against what inference engines do (BF16 tensor cores, FP32 accumulation, split-K
chosen by shape):
\begin{center}\begin{tabular}{lr}\toprule
certified correctly rounded (48 matrices, 79.7~M outputs) & all \\
independent check (Python integers, raw bytes) & 1500/1500 \\
batch 1, 3, 17, 64, 200 vs 4096; decode kernel vs prefill & bit-identical \\
engine outputs not correctly rounded, split-K 1/2/4/8 (expert 0 w1) & 81.5/67.1/55.6/51.2\% \\
engine, split-K 1 vs 4: outputs whose bits differ & 71\% \\
prefill, 4096 tokens, $3072\times3584$: exact / BF16 & 8.0 / 2.9 ms \\
decode, 1 token: exact (packed 4-bit) / FP4 GEMV, FP32 & 0.062 / 0.125 ms \\\bottomrule
\end{tabular}\end{center}
On the real router of the layer ($896\times7168$ BF16, FP32 logits, sigmoid plus learned bias, top 16), the rounding of the
logits alone sends $\sim$1 token in 13\,000 to a different expert (11 in 147\,456 token--split pairs), i.e.\ over K3's 92 MoE
layers $\sim$1 token in 140 runs a different expert somewhere; with exact logits the choice depends on the input alone.
A whole MoE block of layer 5 (62 real experts fetched on demand, the latent projections, RMSNorm and the shared experts;
products correctly rounded, reductions in FP64 in a fixed order) gives the same bits for every token alone or in a batch,
while the engine path changes 49\% of the block's 28\,672 BF16 outputs between split-K 1 and 4.

Where exactness costs nothing: a trillion-parameter model on this laptop is a streaming problem. Per text token K3 reads 25.8~GB
of routed MXFP4 experts and 111.2~GB of dense BF16 weights (attention 72.4, shared experts 24.3; summed from the 96 safetensors headers); we measured 5.1~GB/s from the NVMe drive and 45~GB/s over
PCIe~5~$\times$16, while the exact decode kernel consumes 94~GB/s. A bit-reproducible run therefore takes the same
$\sim$27~s per token as an ordinary one, all of it reading. What this laptop lacks is 1.56~TB of storage (669~GB free).

\paragraph{The largest model with a published data recipe.} Nemotron~3 Ultra~\cite{nemotron3} publishes weights and recipe
and names its pre-training data (560\,524\,578\,816 parameters in the BF16 checkpoint's metadata; 14.8T tokens in 226
datasets), but the bulk of that data (Nemotron-CC v2 and v2.1, the code and SFT pre-training sets) is gated by manual
approval under an agreement that allows internal training only (Hugging Face API, 27 October 2026). Its
routed experts are NVFP4 for weights and activations: E2M1 values, one E4M3 scale per 16, an FP32 global scale. An NVFP4
block of 16 is exactly one \texttt{mma.m16n8k16} INT8 ($|p_b|\le 2304$), the product of the two E4M3 scales is an integer
$\le 225$ times $2^{\varepsilon_W+\varepsilon_A}$, and after a shift of 18 every term is below $2^{47}$: the $K/16\le 320$
terms add in an \texttt{int64} without loss, and the 104-bit product with $g_Wg_A$ is rounded once. On expert~0 of MoE layers
49 and 99 (real weights, 4096 tokens through the real latent projection, up, relu$^2$, down) the tensor-core path, an
integer-loop path, 32 tokens alone, Python rationals from the raw bytes (sampled) and a CPU recomputation of all
29\,360\,128 outputs of layer 49 (BLAS block sums, Python integers) agree in every bit. The hardware's own NVFP4 instruction,
\texttt{mma.m16n8k64.kind::mxf4nvf4.block\_scale.scale\_vec::4X} with FP32 accumulation (what inference engines run), leaves
18--49\% of the outputs not correctly rounded and changes up to 0.81\% of them between split-K 1 and 4; an emulation with
an exact epilogue shows that the FP32 sum alone accounts for 0--0.9\% and the rest is $\alpha=\mathrm{fp32}(g_Wg_A)$. The
native instruction runs at 49--55~TOPS with a plain tiling and the exact path at 20.5--23.2~TOPS with shared-memory tiles
(25.6 without its \texttt{int64} epilogue: the epilogue costs $\sim$20\%, the rest is four \texttt{k16} INT8 MMAs per
\texttt{k64} FP4 one); sustained at $\sim$94~W they give 447--473 and 194--204~GOP per joule. The ceiling is the instruction peak: measured in
registers, \texttt{m16n8k16} INT8 reaches 116~TOPS and the block-scaled FP4 instruction 465~TOPS, so with both kernels tuned
to peak a correctly rounded prefill costs $4\times$; in decode both read the same packed bytes. The exact decode reads
405--500~GB/s of packed weights for one token and 136--184~GB/s for eight; the whole NVFP4 checkpoint (352~GB, from its 113 safetensors headers) fits this laptop's free disk, and a token reads
51.2~GB of it (37.8 dense, 13.4 of routed experts; measured below: 26~s per token from NVMe). Its FP32 router is exactly three BF16
matrices, so its logits are one correctly rounded product with $K=3\cdot 8192$: FP32 accumulation rounds 90--97\% of them
differently, yet only 1 of 16\,384 token-layer routing decisions moves (2 layers, 2 seeds, 4096 tokens).

\paragraph{The whole model, token by token, from disk.} One GEMV covers the four formats of the checkpoint: BF16 (attention,
latent projections, lm\_head), FP32 (router), FP8 E4M3 W8A8 with per-tensor scales and a static input scale (Mamba in and
out, shared experts) and NVFP4 W4A4 (routed experts). Once an activation is quantized as its format prescribes, every
element is an integer times a power of two, so a warp accumulates its row in a 128-bit fixed point anchored at the row's
largest possible term, multiplies by the exact product of the FP32 global scales and rounds once to BF16 (FP32 for router
and logits); terms below the 125-bit window are counted and the rounding is certified when the interval they span rounds to
one value. Against a 2560-bit superaccumulator on the host, 6\,656 outputs in all four formats (narrow and $2^{\pm60}$
ranges, heavy cancellation, exact ties) agree in every bit; a 40-bit window and ties-away rounding are both caught (the tie
test also caught a bug: squeezing the accumulator to 78 bits widened the interval when no bit was lost, so every exact tie
came out uncertified). Everything else (the Mamba-2 recurrence with its FP32 state, the causal convolution, attention
without positional encoding, sigmoid top-22 routing, RMSNorms, residuals) runs on the host in FP64 and is rounded once to
the precision the model keeps there, identically for every path. All 560\,B parameters stream from the laptop's NVMe: one
sequence reads 48~GB per token in 16~s (2.95~GB/s) and takes 26~s, and ``The capital of Kosovo is'' continues
``\,Pristina''. Three prompts (Albanian, English, $123456789\times 987654321$) then ran as nine sequences, each prompt exact
and on the same host code with FP32-accumulating GEMVs (split-K 1 or 4: two configurations of one engine), 45 steps of
57--71~s (70--86~GB read): 1.54 billion correctly rounded outputs, none uncertified. On the exact path's own inputs the
engine rounds 12.6--13.8\% of the BF16 projections and 40.6--47.3\% of the router logits differently, but only 0.003\% of
the FP8 and NVFP4 ones (41\,914 of 1.42 billion): at BF16 output precision FP32 accumulation of the narrow formats is almost
always exact, of the wide ones it is not. Downstream, the engine's top logit drifts up to 1.22 from the exact one (0.17--0.31
on average) while the tokens still agree, the exact top-two margin falls to 0.017, and the streams part: the Albanian one at
the 5th generated token with split-K 4 (an indentation of 8 against 4 spaces) and at the 13th with split-K 1 (``\,ri''
against ``\,pav''; the exact path goes on ``Kosova është vend i ri i pavarur'', the engine ``Kosova është vend i pavarur që
ka deklaruar pavarësinë më 17 shkurt 2008''), the English one at the 29th (``\,and'' against ``\,of''). On the product all
three write 121932631112635269, which is correct. The drivers print a 64-bit FNV-1a digest of every residual stream and
logit vector of the exact path, so anyone with the checkpoint can check a run.

\paragraph{Full access.} Checking an output against the data a model learned from needs the data, not only its recipe.
Measured on the Hugging Face API and the model cards (27 October 2026): Ultra's main sets are gated as above; Jais-3-375B-A23B
(379\,B, Apache-2.0) beats Ultra on most rows of its own card but publishes a 10\% sample of its data (23.8\,B tokens); the
strongest model whose data is wholly public and ungated is IFM's K2-Think-V2 (72.6\,B, Apache-2.0, TxT360 under ODC-BY);
and the one model whose entire training data is indexed for public lookup is OLMo-2-0325-32B-Instruct (all 4.6\,T tokens in
the infini-gram index \texttt{v4\_olmo-2-0325-32b-instruct\_llama}). No fully open model is better than Ultra today. The same
exact GEMV drives both dense models (\texttt{llm/dense.cu}: OLMo~2's post-norm with $q$/$k$ norms; Llama with YaRN, K2's FP32
checkpoint rounded to BF16 as an engine loads it), and every maximal span of an OLMo~2 output is looked up in its own
training data with its count and a source document (\texttt{llm/olmo\_trace.py}): ``The capital of Kosovo is Pristina''
occurs 224 times, one of them a DCLM page from graphicmaps.com.
\fussy

\section{Energy}\label{sec:energy}
Under these kernels the laptop holds $\sim$95~W and the SM clock drops to $\sim$1.8~GHz (2.76~GHz with only the FP64 units busy).
For GEMM at $n=4096$ (\texttt{bench/energy.cu}, power every 200~ms): native FP64 7.4~GFLOPS/W, Ozaki-II at 16 moduli 53.6
($7.2\times$), at 12 moduli 69.5. For a whole solve ($n=4096$, $\kappa=10^{20}$, 20 back-to-back GPU solves after a warm-up; CPU
package power from the RAPL counters, GPU power from \texttt{nvidia-smi}):
\begin{center}\begin{tabular}{lrrrr}\toprule
& time & CPU package & GPU & energy per solve \\\midrule
\texttt{\_\_float128} LU, 24 threads & 55.0 s & 106.5 W & 8.4 W (idle) & 6.3 kJ \\
dd LU on the GPU & 0.238 s & 28.2 W & 68.3 W & \textbf{23 J} \\\bottomrule
\end{tabular}\end{center}
That is $\sim$270$\times$ less energy per solve with both CPU and GPU counted ($360\times$ GPU alone against CPU package alone).
DRAM power is not exposed on this laptop.

\section{Limitations}
The INT8 GEMMs are power-bound on this laptop, so a faster kernel does not translate into time; gains come from fewer
operations (fewer moduli) and less data movement. The cost of the higher precisions grows roughly with the square of the
number of slices, as MPFR's does at these sizes, so the ratio against MPFR does not grow without bound; it is largest
against software emulation of each width. The triple-double solver is panel-bound at $n=4096$. Kernel launches cost
$\sim$7~$\mu$s under WSL2 (0.83~$\mu$s per node in a CUDA graph), which matters for the many small launches of the panel.
All measurements are on one laptop model; the CPU references use the same machine's 24 threads.

\section{Related work}
Ozaki scheme~II~\cite{ozaki2} emulates FP64 GEMM with INT8 and the CRT and reports 7.4--9.8~TFLOPS on an RTX~4090; RIKEN's GEMMul8
is its public code. Follow-up work covers SGEMM and power efficiency on an RTX~5080~\cite{ozaki2sgemm}, FP8 and FP4
variants~\cite{fp8,fp4a,fp4b}, fused kernels reaching 81\% of INT8 peak on Blackwell~\cite{emugemm} and complex
GEMM~\cite{cgemm}. Multimodular exact integer matrix multiplication is classical (FLINT; conversions to and from a residue system written as matrix products by Doliskani, Giorgi, Lebreton and Schost~\cite{dgls}); Close in spirit: Ozaki-II for multiple-precision GEMM on CPU INT8 units (AMX, SVE2-i8mm)~\cite{ozamx}; integer matrix products at various precisions on AMX and IFMA by van der Hoeven and Mezzarobba~\cite{vdhm}; modular products over large prime fields on GPUs by Berthomieu, Graillat, Lesnoff and Mary~\cite{bglm}; EmuGEMM already fuses the modular reduction into the GEMM epilogue~\cite{emugemm}. What we have not found is exact integer GEMM of $10^3$--$10^5$-bit entries on GPU INT8 tensor cores, conversions included, measured against FLINT (to be confirmed by reading these works in full). what is new here is running it, conversions included, on INT8 tensor cores with byte-split residues. Ozaki-II mentions an extension beyond FP64 as possible; we found no implementation or measurement of a
double-double or wider GEMM built this way, nor of linear solves at these condition numbers on this class of hardware.
For low-precision neural networks, exact (Kulisch-style) accumulation was proposed in hardware designs~\cite{johnson18,positron}; batch-invariant kernels~\cite{he25} make LLM inference deterministic across batch sizes on one setup, and say explicitly that they are not hardware or software-version invariant. Correct rounding of every product, as in Section~\ref{sec:llm}, is invariant to both, on commodity INT8 units.
\emph{(References to be re-checked at their source before submission.)}

\section{Conclusion}
On a laptop GPU whose FP64 units are 517 times slower than its INT8 tensor cores, precision becomes a software parameter:
the same INT8 GEMMs yield FP64, double-double, triple-double and wider products, and an ill-conditioned dense solve that needs
quad precision finishes $499\times$ sooner than on the laptop's own CPU, with $\sim$270$\times$ less energy. Exact integer products of up to 131\,063-bit entries come from the same units, $10$--$19\times$ faster than FLINT. Exact determinants and exact rational solutions of int8 systems, where no big number has to enter the GPU, come $85$--$355\times$ faster than FLINT, and with 32--256-bit entries high-order lifting on the same GEMMs keeps the ratio rising with the bits, up to $306\times$. The next steps are
the panel (now the bound of the triple-double solve), CUDA graphs for its launches, and out-of-core sizes beyond the GPU's 12~GB.

\begin{thebibliography}{99}
\bibitem{he25} H.~He and Thinking Machines Lab. Defeating nondeterminism in LLM inference. Blog post, 10 September 2025, \url{https://thinkingmachines.ai/blog/defeating-nondeterminism-in-llm-inference}.
\bibitem{johnson18} J.~Johnson. Rethinking floating point for deep learning. arXiv:1811.01721, 2018.
\bibitem{positron} Z.~Carmichael et al. Deep Positron: a deep neural network using the posit number system. arXiv:1812.01762, 2018.
\bibitem{spex} C.~Lourenco, J.~Chen, E.~Moreno-Centeno, T.~A.~Davis. Algorithm 1021: SPEX Left LU, exactly solving sparse linear
systems via a sparse left-looking integer-preserving LU factorization. \emph{ACM Trans. Math. Softw.} 48(2):1--23, 2022.
doi:10.1145/3519024.
\bibitem{ozaki2} K.~Ozaki, Y.~Uchino, T.~Imamura. Ozaki Scheme II: A GEMM-oriented emulation of floating-point matrix
multiplication using an integer modular technique. arXiv:2504.08009, 2025.
\bibitem{ozaki2sgemm} Y.~Uchino, K.~Ozaki, T.~Imamura. High-performance and power-efficient emulation of matrix
multiplication using INT8 matrix engines. arXiv:2508.03984, 2025.
\bibitem{fp8} Y.~Uchino, K.~Ozaki, T.~Imamura. Double-precision matrix multiplication emulation via Ozaki-II scheme with
FP8 quantization. arXiv:2603.10634, 2026.
\bibitem{fp4a} S.~Hayashi, D.~Mukunoki, T.~Hoshino, T.~Katagiri. DGEMM with Ozaki Scheme I/II on FP4 tensor cores: a
base-13 E2M1 limb representation. arXiv:2608.06812, 2026.
\bibitem{fp4b} S.~Hayashi, D.~Mukunoki, T.~Hoshino, T.~Katagiri. AWE: Adaptive weight encoding for exact integer matrix
products with fewer GEMMs on FP4 tensor cores. arXiv:2609.24519, 2026.
\bibitem{emugemm} D.~Lu, A.~Maeder, M.~Luisier, A.~N.~Ziogas. EmuGEMM: Fused tensor core kernels for precision emulation
in matrix multiplication. arXiv:2606.25453, 2026.
\bibitem{cgemm} Y.~Uchino, Q.~Ma, T.~Imamura, K.~Ozaki, P.~L.~Gutsche. Emulation of complex matrix multiplication based on
the Chinese remainder theorem. arXiv:2512.08321, 2025.
\bibitem{dgls} J.~Doliskani, P.~Giorgi, R.~Lebreton, \'E.~Schost. Simultaneous conversions with the residue number system
using linear algebra. ACM Trans.\ Math.\ Softw.\ 44 (2018). doi:10.1145/3145573.
\bibitem{ozamx} T.~Kouya. Ozaki Scheme II is fast on CPUs too: multiple-precision matrix multiplication on Intel AMX-INT8
and Arm SVE2-i8mm. arXiv:2609.27831, 2026.
\bibitem{vdhm} J.~van der Hoeven, M.~Mezzarobba. Integer matrix multiplication at various precisions using AMX and IFMA.
Preliminary version of June 1, 2026, \url{https://www.texmacs.org/joris/amx/amx.html}.
\bibitem{bglm} J.~Berthomieu, S.~Graillat, D.~Lesnoff, T.~Mary. Multiword matrix multiplication over large finite fields
in floating-point arithmetic. arXiv:2601.07508, 2026.
\bibitem{ogita} T.~Ogita. Accurate and verified numerical computation of the matrix determinant. Int.\ J.\ Reliability and
Safety 6(1/2/3) (2012) 242. doi:10.1504/IJRS.2012.044287.
\bibitem{pan} A.~Urba\'nska. Towards an exact adaptive algorithm for the determinant of a rational matrix.
arXiv:0706.0014, 2007.
\bibitem{stor05} A.~Storjohann. The shifted number system for fast linear algebra on integer matrices. J.\ Complexity 21(4):609--650, 2005. doi:10.1016/j.jco.2005.04.002
\bibitem{ps} C.~Pauderis, A.~Storjohann. Deterministic unimodularity certification. Proc.\ ISSAC 2012, 281--288.
doi:10.1145/2442829.2442870.
\bibitem{nemotron3} NVIDIA. NVIDIA-Nemotron-3-Ultra-550B-A55B (model card; BF16 and NVFP4 checkpoints). Hugging Face, 3 June 2026,
\url{https://huggingface.co/nvidia/NVIDIA-Nemotron-3-Ultra-550B-A55B-NVFP4}.
\end{thebibliography}

\end{document}
